\section{Discussion}
\label{sec:discussion}

\subsection{The Constraint-Quality Tradeoff}
\label{sec:tradeoff}

Our central finding is a \emph{constraint-quality tradeoff}: as more structural constraints are imposed, the model follows them faithfully but produces less creative, less linguistically polished output.
Qualitative examination reveals the mechanism.
Baseline stories use vivid imagery, sensory language, and unexpected turns---the model invests its generation capacity in expressive writing.
Component-controlled stories read more like competent plot summaries, hitting all specification points but with flatter prose.
The model appears to allocate its generation ``bandwidth'' between two competing objectives: (1) satisfying constraints and (2) producing expressive writing.
More constraints leave less room for creative expression.

This tradeoff has a natural interpretation in terms of conditional generation difficulty.
The unconstrained model samples from $p(\text{story} \mid \text{theme})$, a broad distribution with many high-quality modes.
The constrained model samples from $p(\text{story} \mid \text{setting}, \text{character}, \text{conflict}, \text{theme}, \text{arc}, \text{tone})$, a much narrower distribution where the model must navigate a smaller space of acceptable outputs.
The constraint-quality tradeoff reflects the difficulty of finding outputs that simultaneously satisfy all constraints \emph{and} exhibit high linguistic quality.

\subsection{Implications for the Arbitrage Hypothesis}

Our results refute the quality arbitrage hypothesis (H4) for single-pass generation: models that write mediocre stories do \emph{not} write better stories when given explicit structural specifications.
However, the arbitrage may exist in a \emph{two-stage} pipeline.
If the component specification stage provides the structural backbone and a subsequent unconstrained refinement stage polishes the prose, the combined system could achieve both controllability and quality.
This echoes findings from DSR~\cite{dsr2025}, where separating narrative creation from format conversion yielded 82.7\% of human-level quality, and from Agents' Room~\cite{huot2024agentsroom}, where multi-agent decomposition outperformed single-model approaches.

\subsection{Implications for Synthetic Data Generation}

Despite the quality penalty, our pipeline has clear value for synthetic data generation.
The extracted component taxonomy provides a structured vocabulary for creating diverse training specifications.
The near-perfect controllability (4.96/5.0) means the generated stories reliably exhibit the specified characteristics, making them suitable for training smaller models to follow narrative specifications via distillation.
The quality penalty may be acceptable for training data---the structural diversity of the data matters more than the polish of individual examples, especially if the student model is separately trained on quality.

\subsection{Limitations}
\label{sec:limitations}

\para{Same-model evaluation.}
We use \gptfour for both generation and evaluation, creating potential self-preference bias.
While this bias affects all conditions equally and should not change relative rankings, the absolute scores may be inflated.
GPT-4-as-judge shows moderate correlation with human judgments ($\rho \approx 0.4$--$0.6$ per the literature~\cite{huot2024agentsroom,xie2024weaver}), and the high absolute scores (rarely below 3/5) suggest the judge may lack discriminative power.

\para{Short story format.}
\rocstories consists of five-sentence stories averaging 43 words.
This format may not reveal benefits of component-based control that emerge at longer scales, where planning and structural coherence become more important~\cite{yang2023doc,huot2024agentsroom}.
The ceiling effect on coherence (all conditions achieve 5.0/5.0) further suggests that this format is too easy to test the full value of structural control.

\para{No fine-tuning.}
We test zero-shot component-controlled generation.
The \skillmix paradigm's full potential requires training on synthetic component data, which we did not test due to scope constraints.
The observed quality penalty might diminish or reverse after fine-tuning.

\para{Single model.}
All experiments use \gptfour.
Results may differ with other models, especially smaller models where controllability is harder to achieve and where the benefits of structural guidance may be larger.

\para{Automatic evaluation only.}
No human evaluation was conducted.
Human judges might disagree with the \gptfour judge's relative rankings, particularly on subjective dimensions like creativity and language quality.

\para{Limited taxonomy.}
The narrative technique dimension collapsed to a single value (linear) due to \rocstories' short format, reducing the effective component space from seven to six dimensions.
